% ============================================
% RÉSULTATS DE VALIDATION - COMPARAISON DDQN
% ============================================

\section{Résultats de Validation}

\subsection{Comparaison Approche DDQN vs Approche Standard}

Les résultats de validation montrent une amélioration significative de la stratégie basée sur l'agent DDQN comparée à une approche standard avec puissance fixe.

\subsubsection{Packet Delivery Ratio (PDR)}

L'agent DDQN atteint un PDR moyen de \textbf{95.92\%} sur l'ensemble des distances testées, contre 92.73\% pour l'approche standard. Cela représente une \textbf{amélioration moyenne de +3.20\%}.

\vspace{0.3cm}

\noindent \textbf{Analyse par distance:}
\begin{itemize}
    \item À courte portée (139-331m): PDR quasi identique (\textasciitilde 99-100\%)
    \item À moyenne portée (775-1977m): DDQN supérieur ou égal dans 8 cas sur 9
    \item À longue portée (2000-4000m): \textbf{Réduction de perte de 5-15\%} 
    \item Cas critique (4000m): DDQN atteint 92.60\% vs 77.88\% en standard (\textbf{+14.73\%})
\end{itemize}

\begin{figure}[h]
    \centering
    \includegraphics[width=0.95\textwidth]{images/results/pdr_comparison.png}
    \caption{Comparaison du Packet Delivery Ratio entre l'approche standard et l'agent DDQN sur 13 distances différentes. DDQN montre une supériorité particulièrement marquée à très longue portée.}
    \label{fig:pdr_comparison}
\end{figure}

\begin{figure}[h]
    \centering
    \includegraphics[width=0.95\textwidth]{images/results/pdr_gain_ddqn.png}
    \caption{Gain en PDR apporté par l'agent DDQN par rapport à l'approche standard. Les valeurs positives (vert) indiquent une supériorité de DDQN. Le gain atteint +14.73\% à 4000m.}
    \label{fig:pdr_gain}
\end{figure}

\subsubsection{Efficacité Énergétique}

La consommation d'énergie moyenne est réduite de \textbf{5.67\%} avec l'approche DDQN:
\begin{itemize}
    \item DDQN: \textbf{72.53 J} en moyenne
    \item Standard: 76.90 J en moyenne
    \item Économie cumulée sur 13 distances: \textbf{56.74 J}
\end{itemize}

Cependant, à très longue portée (> 1194m), l'agent augmente légèrement la puissance pour maintenir la fiabilité, d'où une consommation légèrement supérieure. Cette stratégie d'échange \"fiabilité-énergie\" est optimale pour maximiser simultanément PDR et efficacité énergétique.
\begin{figure}[h]
    \centering
    \includegraphics[width=0.95\textwidth]{images/results/energy_comparison.png}
    \caption{Consommation d'énergie par approche et par distance. À courte-moyenne portée, DDQN économise l'énergie grâce à une puissance réduite. À longue portée, la puissance augmente pour maintenir la fiabilité, justifiant une augmentation contrôlée de la consommation.}
    \label{fig:energy_comparison}
\end{figure}

\begin{figure}[h]
    \centering
    \includegraphics[width=0.95\textwidth]{images/results/energy_reduction.png}
    \caption{Réduction relative de consommation d'énergie avec DDQN. Les valeurs positives (bleu) indiquent une efficacité supérieure de DDQN. La réduction atteint jusqu'à 62\% à courte portée.}
    \label{fig:energy_reduction}
\end{figure}
\subsubsection{Stratégie d'Adaptation de Puissance}

L'agent DDQN ajuste la puissance de transmission de manière intelligente:
\begin{itemize}
    \item Puissance moyenne DDQN: \textbf{10.24 dBm}
    \item Puissance standard: 14.00 dBm
    \item Réduction moyenne: \textbf{3.76 dBm}
\end{itemize}

À courte-moyenne portée, l'agent réduit la puissance (3-7 dBm) pour économiser l'énergie, tandis qu'à très longue portée (4000m), il l'augmente à 14 dBm pour maintenir la connectivité.

\begin{figure}[h]
    \centering
    \includegraphics[width=0.95\textwidth]{images/results/power_adaptation.png}
    \caption{Stratégie adaptative de puissance de l'agent DDQN comparée à l'approche standard fixe. Le comportement intelligent de l'agent est clairement visible: puissance réduite à courte portée pour économiser, puis augmentation progressive avec la distance pour maintenir la fiabilité.}
    \label{fig:power_adaptation}
\end{figure}

\subsection{Validation du Simulateur LR-FHSS}

\subsubsection{Couverture et Portée}

Le simulateur a validé une portée fiable (PDR > 85\%) jusqu'à \textbf{4000m} avec:
\begin{itemize}
    \item PDR global moyen: 93.20\%
    \item PDR minimum (4000m, DR11): 63.78\%
    \item Tendance cohérente de dégradation avec la distance
\end{itemize}

\begin{figure}[h]
    \centering
    \includegraphics[width=0.95\textwidth]{images/results/pdr_vs_distance.png}
    \caption{Dégradation du PDR en fonction de la distance pour l'ensemble des configurations DR/CR. La courbe montre une tendance de dégradation naturelle avec la distance, validant la réalisme du simulateur. La zone verte indique la plage de fiabilité (PDR > 85\%).}
    \label{fig:pdr_distance}
\end{figure}

\subsubsection{Performance par Data Rate}

\begin{table}[h]
\centering
\small
\begin{tabular}{|c|c|c|c|c|}
\hline
\textbf{Data Rate} & \textbf{PDR Moyen} & \textbf{Plage} & \textbf{RSSI Moyen} & \textbf{BER Moyen} \\
\hline
DR8 (CR 1/3) & 96.99\% & 90.27-100\% & -114.94 dB & 9.81e-03 \\
\hline
DR9 (CR 2/3) & 93.28\% & 80.44-100\% & -114.58 dB & 1.45e-02 \\
\hline
DR10 (CR 1/3) & 94.49\% & 81.88-100\% & -114.65 dB & 1.79e-02 \\
\hline
DR11 (CR 2/3) & 88.02\% & 63.78-100\% & -114.34 dB & 2.71e-02 \\
\hline
\end{tabular}
\caption{Performance en fonction du Data Rate}
\label{tab:dr_performance}
\end{table}

Les résultats montrent que:
\begin{itemize}
    \item \textbf{DR8 (CR 1/3)} offre les meilleures performances globales avec un codage robuste
    \item \textbf{DR9 et DR10} maintiennent une couverture acceptable (PDR > 93\%)
    \item \textbf{DR11 (CR 2/3)} est plus sensible aux dégradations à longue portée due au codage moins redondant
\end{itemize}

\begin{figure}[h]
    \centering
    \includegraphics[width=0.95\textwidth]{images/results/pdr_by_datarate.png}
    \caption{Performance du PDR par Data Rate avec leur plage de variation (min-max). DR8 affiche la meilleure performance globale, tandis que DR11 montre la plus grande variabilité, reflétant sa sensibilité aux conditions de propagation.}
    \label{fig:pdr_datarate}
\end{figure}

\subsubsection{Dégradation du Signal}

L'analyse du RSSI (Received Signal Strength Indicator) et du SNR (Signal-to-Noise Ratio) montre une tendance cohérente avec la théorie:

\begin{table}[h]
\centering
\small
\begin{tabular}{|c|c|c|c|}
\hline
\textbf{Plage de Distance} & \textbf{PDR Moyen} & \textbf{RSSI Moyen} & \textbf{SNR Moyen} \\
\hline
Courte (139-331m) & 100\% & -83.7 à -95.4 dB & 19-31 dB \\
\hline
Moyenne (775-1300m) & 97.7\% & -107 à -115 dB & -1 à 7 dB \\
\hline
Longue (1753-4000m) & 86.4\% & -122 à -130 dB & -12 à -8 dB \\
\hline
\end{tabular}
\caption{Performance par plage de distance}
\label{tab:distance_performance}
\end{table}

\subsection{Conclusions de la Validation}

Les résultats de validation confirment:

\begin{enumerate}
    \item \textbf{Efficacité de l'agent DDQN}: L'approche par apprentissage par renforcement améliore le PDR de 3-5\% en moyenne, et jusqu'à 15\% à très longue portée.
    
    \item \textbf{Pertinence de la simulation}: Les tendances observées (dégradation avec la distance, performance par DR) sont cohérentes avec les études de référence et les modèles théoriques.
    
    \item \textbf{Équilibre fiabilité-efficacité}: L'agent DDQN réalise un compromis optimal en réduisant la puissance à courte portée (économie d'énergie) tout en l'augmentant à longue portée (maintien de la fiabilité).
    
    \item \textbf{Applicabilité pratique}: La portée fiable jusqu'à 4 km confirme la viabilité du LR-FHSS pour les applications IoT longue portée.
\end{enumerate}
